% ============================================================
%  DOCUMENTO DE PLANIFICACIÓN XP --- VitaPrevent
%  Metodología Extreme Programming: Iteraciones reales + plan
%  Asignatura: Usabilidad y Accesibilidad (ISWD732) – GR3SW
%  EPN – Facultad de Ingeniería de Sistemas – Ingeniería de Software
%  Compilar con: pdflatex metodologia_xp_vitaprevent.tex (2 veces)
% ============================================================
\documentclass[12pt, a4paper]{article}

\usepackage[utf8]{inputenc}
\usepackage[T1]{fontenc}
\usepackage[spanish, es-noshorthands]{babel}
\usepackage[top=2.5cm, bottom=2.5cm, left=3cm, right=2.5cm, headheight=14pt]{geometry}

\usepackage[dvipsnames, table]{xcolor}
\definecolor{navydeep}{HTML}{0D1B2A}
\definecolor{royalblue}{HTML}{1565C0}
\definecolor{lightblue}{HTML}{BBDEFB}
\definecolor{successgreen}{HTML}{16A34A}
\definecolor{warningyellow}{HTML}{D97706}
\definecolor{errored}{HTML}{DC2626}
\definecolor{graymid}{HTML}{6B7280}
\definecolor{graylight}{HTML}{F3F4F6}
\definecolor{violetxp}{HTML}{6D28D9}
\definecolor{tealxp}{HTML}{0F766E}

\usepackage{graphicx}
\usepackage{float}
\usepackage{tikz}
\usetikzlibrary{shapes, arrows.meta, positioning, decorations.pathreplacing, calc}
\usepackage{pgfplots}
\pgfplotsset{compat=1.18}

\usepackage{booktabs}
\usepackage{longtable}
\usepackage{array}
\usepackage{multirow}
\usepackage{makecell}
\renewcommand{\arraystretch}{1.5}

\usepackage{listings}
\lstdefinestyle{gitlog}{
  backgroundcolor=\color{navydeep!6},
  basicstyle=\ttfamily\footnotesize,
  breaklines=true, frame=single, framesep=4pt,
  rulecolor=\color{royalblue!30},
  xleftmargin=6pt, xrightmargin=4pt,
  tabsize=2, showstringspaces=false,
  literate=
    {á}{{\'a}}1 {é}{{\'e}}1 {í}{{\'i}}1 {ó}{{\'o}}1 {ú}{{\'u}}1
    {Á}{{\'A}}1 {É}{{\'E}}1 {Í}{{\'I}}1 {Ó}{{\'O}}1 {Ú}{{\'U}}1
    {ñ}{{\~n}}1 {Ñ}{{\~N}}1,
}
\lstset{style=gitlog}

\usepackage[
  colorlinks=true, linkcolor=royalblue, urlcolor=royalblue, citecolor=royalblue,
  unicode=true, pdfencoding=auto,
  pdftitle={VitaPrevent - Metodología XP - Planificación de Iteraciones},
  pdfauthor={Erick Costa, Jonathan Tipán, Javier Quilumba},
  pdflang={es-EC}
]{hyperref}
\usepackage{bookmark}
\usepackage{amssymb}
\newcommand{\cmark}{\textcolor{successgreen}{$\checkmark$}}
\newcommand{\wmark}{\textcolor{warningyellow}{$\triangleright$}}
\newcommand{\pmark}{\textcolor{royalblue}{$\bullet$}}

\usepackage{fancyhdr}
\pagestyle{fancy}
\fancyhf{}
\fancyhead[L]{\small\color{graymid}\textit{VitaPrevent --- Metodología XP --- ISWD732 GR3SW}}
\fancyhead[R]{\small\color{graymid}\thepage}
\fancyfoot[C]{\small\color{graymid}\textit{Escuela Politécnica Nacional --- 2026}}
\renewcommand{\headrulewidth}{0.4pt}
\renewcommand{\footrulewidth}{0.4pt}

\usepackage{titlesec}
\titleformat{\section}
  {\normalfont\large\bfseries\color{navydeep}}
  {\thesection.}{0.6em}{}[{\color{royalblue}\titlerule[0.8pt]}]
\titleformat{\subsection}
  {\normalfont\normalsize\bfseries\color{royalblue}}
  {\thesubsection}{0.6em}{}
\titleformat{\subsubsection}
  {\normalfont\small\bfseries\color{violetxp}}
  {\thesubsubsection}{0.6em}{}

\usepackage[most]{tcolorbox}
\tcbuselibrary{skins, breakable}
\newtcolorbox{storybox}[2]{enhanced, breakable,
  colback=#2!5, colframe=#2,
  fonttitle=\bfseries\small\color{white}, coltitle=white,
  title={\textsf{#1}},
  arc=4pt, boxrule=0.8pt, left=8pt, right=8pt, top=6pt, bottom=6pt}
\newtcolorbox{iterbox}[2]{enhanced, breakable,
  colback=#2!4, colframe=#2,
  fonttitle=\bfseries\normalsize\color{white}, coltitle=white,
  title={#1}, arc=6pt, boxrule=1pt,
  left=10pt, right=10pt, top=8pt, bottom=8pt}
\newtcolorbox{metabox}[1][]{enhanced,
  colback=graylight, colframe=graymid!60,
  fonttitle=\bfseries\small, title={#1},
  arc=4pt, boxrule=0.6pt, left=8pt, right=8pt, top=6pt, bottom=6pt}

\usepackage{enumitem}
\setlist[itemize]{leftmargin=1.4em, itemsep=3pt}
\setlist[enumerate]{leftmargin=1.4em, itemsep=3pt}

\usepackage{setspace}
\setstretch{1.25}
\usepackage{parskip}
\setlength{\parindent}{0pt}
\setlength{\parskip}{6pt}
\usepackage{caption}
\captionsetup{font=small, labelfont={bf, color=navydeep}}

% ============================================================
\begin{document}
\pagenumbering{gobble}

% ------------------------------------------------------------
%  PORTADA
% ------------------------------------------------------------
\begin{titlepage}
\centering
\vspace*{0.5cm}

\begin{tikzpicture}
  \fill[royalblue, rounded corners=8pt] (0,0) rectangle (2.0,2.0);
  \fill[white] (0.85,0.25) rectangle (1.15,1.75);
  \fill[white] (0.25,0.85) rectangle (1.75,1.15);
  \fill[royalblue] (0.9,0.9) circle (0.22);
\end{tikzpicture}

\vspace{0.5cm}
{\LARGE\bfseries\color{navydeep} ESCUELA POLITÉCNICA NACIONAL}\\[4pt]
{\large\color{navydeep} FACULTAD DE INGENIERÍA DE SISTEMAS}\\[2pt]
{\large\color{navydeep} INGENIERÍA DE SOFTWARE}

\vspace{1.0cm}
{\color{violetxp}\rule{0.85\linewidth}{1.5pt}}

\vspace{0.6cm}
{\large\bfseries\color{navydeep} TAREA 06 --- DOCUMENTO COMPLEMENTARIO}\\[6pt]
{\Large\bfseries\color{navydeep}
  PLANIFICACIÓN Y GESTIÓN DEL PROYECTO\\
  MEDIANTE METODOLOGÍA\\[4pt]
  \textcolor{violetxp}{EXTREME PROGRAMMING (XP)}}

\vspace{0.8cm}
{\color{violetxp}\rule{0.85\linewidth}{0.8pt}}

\vspace{0.8cm}
\begin{tabular}{@{}rl@{}}
  \textbf{\color{navydeep}Asignatura:} & Usabilidad y Accesibilidad (ISWD732) \\[4pt]
  \textbf{\color{navydeep}Curso:}      & GR3SW \\[8pt]
\end{tabular}

\vspace{0.4cm}
\begin{center}
{\normalsize\bfseries\color{navydeep} Proyecto:}\\[4pt]
{\large\bfseries\color{royalblue}
  VITAPREVENT: PLATAFORMA WEB ACCESIBLE\\
  DE SERVICIOS PÚBLICOS DE SALUD}
\end{center}

\vspace{0.8cm}
\begin{tabular}{@{}rl@{}}
  \textbf{\color{navydeep}Grupo 6:} & Javier Quilumba \\
  & Jonathan Tipán \\
  & Erick Costa \\[8pt]
  \textbf{\color{navydeep}Docente:} & Ing. Pablo Andrés Saá P. \\[4pt]
  \textbf{\color{navydeep}Fecha:}   & 28/06/2026 \\
\end{tabular}

\vspace{1.0cm}
{\color{violetxp}\rule{0.85\linewidth}{1.5pt}}
\end{titlepage}

% ------------------------------------------------------------
%  TABLA DE CONTENIDOS
% ------------------------------------------------------------
\newpage
\pagenumbering{roman}
\tableofcontents

\newpage
\pagenumbering{arabic}

% ============================================================
\section{Introducción}

El presente documento describe la planificación, ejecución y gestión del proyecto VitaPrevent empleando la metodología \textbf{Extreme Programming (XP)}, seleccionada por su adaptabilidad a equipos pequeños, su énfasis en la calidad del código y su capacidad para incorporar cambios de requerimientos durante el desarrollo, algo especialmente relevante en un proyecto académico con restricciones de tiempo.

VitaPrevent es una plataforma web accesible de servicios públicos de salud del Ecuador, construida con React 19 + Vite 6 + Firebase. El desarrollo real se ejecutó en \textbf{cuatro iteraciones} comprendidas entre el 28/06/2026 y el 29/06/2026. A partir de esa base funcional, el Sprint 2 (30/06/2026 -- 13/07/2026) se encuentra \textbf{en ejecución}. Al 02/07/2026 ya se han completado las historias HU-11, HU-14 y parte de HU-13, y se han resuelto cinco correcciones de calidad relacionadas con accesibilidad de imágenes y renderizado semántico de contenido.

Este documento incluye:
\begin{itemize}
  \item El conjunto de historias de usuario que guió el desarrollo.
  \item Los resultados reales de la Iteración 1 (Sprint 1), con evidencia del historial de commits.
  \item El plan de la Iteración 2 (Sprint 2) con tareas, responsables y criterios de aceptación.
  \item Las métricas de velocidad del equipo y el gráfico de avance (burndown).
  \item Las prácticas XP aplicadas a lo largo del proyecto.
\end{itemize}

% ============================================================
\section{Fundamentos de la Metodología XP}

\subsection{¿Por qué XP?}

Extreme Programming fue elegida sobre metodologías más pesadas (RUP, cascada) porque:
\begin{itemize}
  \item El equipo es pequeño (3 personas), exactamente el escenario para el que XP fue diseñado.
  \item Los requerimientos de accesibilidad evolucionaron durante el desarrollo (reenfoque de salud preventiva genérica a servicios públicos de Ecuador).
  \item XP prioriza \textbf{entregas frecuentes} (cada iteración produce software funcional y desplegado), compatible con el cronograma académico.
  \item Las prácticas de \textbf{integración continua} y \textbf{código limpio} son directamente aplicables con el stack React + GitHub Actions.
\end{itemize}

\subsection{Valores XP aplicados al proyecto}

\begin{table}[H]
\centering
\caption{Valores XP y su aplicación en VitaPrevent}
\small
\begin{tabular}{@{}>{\bfseries\color{violetxp}}p{3cm} p{11.5cm}@{}}
\toprule
\rowcolor{violetxp!10}
\textbf{Valor XP} & \textbf{Aplicación concreta en VitaPrevent} \\
\midrule
Comunicación &
  Repositorio GitHub compartido con commits descriptivos. Todos los integrantes conocen el estado completo del proyecto. Pull del código antes de cada sesión. \\
Simplicidad &
  Sin frameworks visuales pesados (no Bootstrap, no Material UI). Componentes con una sola responsabilidad. CSS con variables centralizadas. \\
Retroalimentación &
  axe DevTools detecta violaciones WCAG en tiempo real. \texttt{eslint-plugin-jsx-a11y} bloquea builds con errores de accesibilidad. Pruebas con NVDA al finalizar cada iteración. \\
Coraje &
  Refactorización del módulo de rutas en Iteración 2 (de rutas planas a rutas anidadas con ProtectedRoute). Reenfoque completo del tema en Iteración 4. \\
Respeto &
  Commits alternados entre los tres integrantes. Nadie sobreescribe trabajo ajeno sin coordinación. \\
\bottomrule
\end{tabular}
\end{table}

\subsection{Prácticas XP implementadas}

\begin{enumerate}
  \item \textbf{Juego de planificación:} Historias de usuario estimadas y priorizadas al inicio de cada iteración.
  \item \textbf{Pequeñas entregas:} El sitio se despliega en Firebase Hosting al finalizar cada iteración (URL pública funcional).
  \item \textbf{Diseño simple:} Solo se implementa lo que la historia de usuario requiere.
  \item \textbf{Refactorización continua:} Módulos de rutas, servicios Firebase y componentes fueron refactorizados cuando el crecimiento lo exigió.
  \item \textbf{Integración continua:} GitHub Actions ejecuta build y deploy automáticamente en cada push a \texttt{main}.
  \item \textbf{Estándares de código:} ESLint + Prettier + \texttt{jsx-a11y}. El build falla si hay violaciones.
  \item \textbf{Propiedad colectiva del código:} Todos los integrantes modifican cualquier parte del proyecto.
  \item \textbf{Semanas de 40 horas:} Dado el contexto académico, se trabajó en sesiones intensivas y se alternaron los commits para distribuir la carga.
\end{enumerate}

% ============================================================
\section{Equipo y Roles}

\begin{table}[H]
\centering
\caption{Equipo XP de VitaPrevent}
\begin{tabular}{@{}>{\bfseries}p{3.5cm} p{4.5cm} p{6cm}@{}}
\toprule
\rowcolor{violetxp!10}
\textbf{Integrante} & \textbf{Rol XP} & \textbf{Contribuciones principales} \\
\midrule
Erick Costa &
  Tracker / Arquitecto &
  Arquitectura base (Vite + React + Firebase), layout y componentes UI, correcciones WCAG 2.4.3, EC\_SALUD data. \\[4pt]
Jonathan Tipán &
  Programador / Integrador &
  Módulos de contenido, panel admin rutas anidadas, CI/CD GitHub Actions, cubo 3D rAF, contraste dark mode. \\[4pt]
Javier Quilumba &
  Coach / Redactor &
  Panel admin CRUD, Biblioteca y Noticias, Firebase Hosting config, alt/title imágenes, documentación. \\
\bottomrule
\end{tabular}
\end{table}

% ============================================================
\section{Historias de Usuario}

Las historias de usuario fueron definidas al inicio del proyecto. Cada historia sigue el formato: \textit{Como [rol], quiero [funcionalidad], para [beneficio]}.

\subsection{Épica 1 --- Estructura y Navegación}

\begin{storybox}{HU-01 \quad Navegación accesible por teclado}{royalblue}
\textbf{Como} usuario que navega solo con teclado,
\textbf{quiero} poder acceder a todas las secciones del sitio usando Tab/Shift+Tab/Enter,
\textbf{para} no depender del mouse.\\[4pt]
\textbf{Criterios de aceptación:}
\begin{itemize}
  \item SkipLink es el primer elemento enfocable en todas las páginas.
  \item Orden de foco: SkipLink $\to$ Navbar $\to$ contenido principal.
  \item Al navegar entre rutas SPA, el foco se mueve al contenido principal (no al inicio del DOM).
  \item Modales capturan el foco y lo devuelven al cerrar.
\end{itemize}
\textbf{Puntos:} 5 \quad \textbf{Estado:} \cmark~Completado (Iteración 4)
\end{storybox}

\begin{storybox}{HU-02 \quad Título de página por ruta}{royalblue}
\textbf{Como} usuario de lector de pantalla,
\textbf{quiero} que cada página tenga un título descriptivo único,
\textbf{para} saber en qué sección me encuentro al navegar.\\[4pt]
\textbf{Criterios de aceptación:}
\begin{itemize}
  \item \texttt{document.title} se actualiza en cada cambio de ruta con formato ``Sección | VitaPrevent''.
  \item La página de inicio usa ``VitaPrevent --- Servicios públicos de salud''.
\end{itemize}
\textbf{Puntos:} 2 \quad \textbf{Estado:} \cmark~Completado (Iteración 1)
\end{storybox}

\subsection{Épica 2 --- Contenido Accesible}

\begin{storybox}{HU-03 \quad Imágenes con texto alternativo}{successgreen}
\textbf{Como} usuario con discapacidad visual,
\textbf{quiero} que todas las imágenes del sitio tengan una descripción,
\textbf{para} que el lector de pantalla me diga qué muestra cada imagen.\\[4pt]
\textbf{Criterios de aceptación:}
\begin{itemize}
  \item Todas las imágenes informativas tienen atributo \texttt{alt} con descripción real (no vacío, no genérico).
  \item Al pasar el mouse sobre una imagen aparece un tooltip con la descripción (atributo \texttt{title}).
  \item Imágenes en contenido Markdown usan un renderer personalizado que preserva alt/title.
  \item Imágenes decorativas usan \texttt{aria-hidden="true"}.
\end{itemize}
\textbf{Puntos:} 3 \quad \textbf{Estado:} \cmark~Completado (Iteración 4)
\end{storybox}

\begin{storybox}{HU-04 \quad Contenido con datos oficiales de salud}{successgreen}
\textbf{Como} ciudadano ecuatoriano,
\textbf{quiero} ver estadísticas de salud de fuentes oficiales,
\textbf{para} conocer la situación real del sistema de salud del país.\\[4pt]
\textbf{Criterios de aceptación:}
\begin{itemize}
  \item La página de inicio muestra al menos 4 estadísticas con fuente y año citados (MSP, IESS, ECU 911).
  \item Cada estadística incluye una barra de progreso visual y texto legible.
\end{itemize}
\textbf{Puntos:} 3 \quad \textbf{Estado:} \cmark~Completado (Iteración 4)
\end{storybox}

\subsection{Épica 3 --- Servicios Públicos de Salud}

\begin{storybox}{HU-05 \quad Módulos de servicios públicos}{tealxp}
\textbf{Como} usuario del sitio,
\textbf{quiero} acceder a información organizada por tipo de servicio público de salud,
\textbf{para} encontrar lo que necesito sin perderme en el sitio.\\[4pt]
\textbf{Criterios de aceptación:}
\begin{itemize}
  \item Existen cuatro módulos: Atención Primaria, Vacunación, Salud Mental y Emergencias.
  \item Cada módulo tiene artículos filtrados por categoría, FAQs en Accordion y enlace de regreso.
  \item Los artículos se cargan desde Firestore de forma dinámica con paginación.
\end{itemize}
\textbf{Puntos:} 8 \quad \textbf{Estado:} \cmark~Completado (Iteraciones 2 y 4)
\end{storybox}

\begin{storybox}{HU-06 \quad Noticias de salud pública}{tealxp}
\textbf{Como} usuario interesado en salud,
\textbf{quiero} ver noticias actualizadas sobre salud pública en Ecuador,
\textbf{para} estar informado sobre campañas y alertas del MSP.\\[4pt]
\textbf{Criterios de aceptación:}
\begin{itemize}
  \item Feed de noticias paginado con cursor de Firestore.
  \item Cada noticia muestra imagen accesible (alt descriptivo), categoría y fecha.
  \item La página de detalle de noticia tiene navegación de regreso y contenido Markdown renderizado.
\end{itemize}
\textbf{Puntos:} 5 \quad \textbf{Estado:} \cmark~Completado (Iteraciones 2 y 3)
\end{storybox}

\subsection{Épica 4 --- Formularios y Accesibilidad}

\begin{storybox}{HU-07 \quad Formulario de contacto accesible}{warningyellow}
\textbf{Como} ciudadano con discapacidad,
\textbf{quiero} enviar consultas usando solo teclado y con ayuda del lector de pantalla,
\textbf{para} comunicarme con los responsables del sitio sin barreras.\\[4pt]
\textbf{Criterios de aceptación:}
\begin{itemize}
  \item Todos los campos tienen \texttt{<label>} visible vinculada con \texttt{htmlFor}.
  \item Errores de validación: texto descriptivo + \texttt{role="alert"} + \texttt{aria-describedby}.
  \item Confirmación de envío anunciada con \texttt{aria-live="assertive"}.
  \item Validación no depende solo del color.
\end{itemize}
\textbf{Puntos:} 5 \quad \textbf{Estado:} \cmark~Completado (Iteración 1)
\end{storybox}

\begin{storybox}{HU-08 \quad Panel administrativo seguro}{warningyellow}
\textbf{Como} administrador del sitio,
\textbf{quiero} gestionar artículos, noticias y recursos desde un panel protegido,
\textbf{para} actualizar el contenido sin acceder directamente a Firebase.\\[4pt]
\textbf{Criterios de aceptación:}
\begin{itemize}
  \item Autenticación con Firebase Auth (email/password). Ruta protegida con ProtectedRoute.
  \item CRUD completo de artículos, noticias y recursos.
  \item Estadísticas en tiempo real desde Firestore.
  \item Notificaciones Toast al crear/editar/eliminar.
\end{itemize}
\textbf{Puntos:} 8 \quad \textbf{Estado:} \cmark~Completado (Iteraciones 2 y 3)
\end{storybox}

\subsection{Historia completada en Iteración 1.4 --- complemento}

\begin{storybox}{HU-03b \quad Imágenes únicas por sección (extensión WCAG 1.1.1)}{successgreen}
\textbf{Como} usuario con discapacidad visual,
\textbf{quiero} que cada artículo tenga una imagen visualmente distinta a los demás de la misma sección,
\textbf{para} que el lector de pantalla no anuncie el mismo \texttt{alt} dos veces para artículos diferentes.\\[4pt]
\textbf{Criterios de aceptación:}
\begin{itemize}
  \item Ningún artículo en la misma sección comparte ID de foto con otro artículo del mismo módulo.
  \item El texto alternativo describe el tema específico del artículo, no un texto genérico.
  \item El servicio \texttt{cleanup.service.js} asigna automáticamente imágenes diferenciadas al ejecutar limpieza.
  \item Para artículos del mismo tema con títulos distintos (ej.~dos artículos sobre dieta mediterránea), se aplica sub-routing por palabras clave del título.
\end{itemize}
\textbf{Puntos:} 3 \quad \textbf{Estado:} \cmark~Completado (Sprint 2, 01/07/2026)
\end{storybox}

\begin{storybox}{HU-11b \quad Renderizado semántico del contenido WYSIWYG (WCAG 1.3.1)}{successgreen}
\textbf{Como} usuario de lector de pantalla,
\textbf{quiero} que el cuerpo de los artículos y noticias se renderice como HTML semántico,
\textbf{para} que el lector de pantalla interprete correctamente encabezados, párrafos y énfasis del contenido.\\[4pt]
\textbf{Criterios de aceptación:}
\begin{itemize}
  \item Artículos con contenido HTML del editor WYSIWYG se renderizan con \texttt{dangerouslySetInnerHTML} (no como texto plano con etiquetas visibles).
  \item Artículos con contenido Markdown se renderizan con \texttt{ReactMarkdown}.
  \item La detección es automática: si el contenido empieza con \texttt{<}, se usa HTML; si no, Markdown.
  \item Aplicado en \texttt{ArticleDetail.jsx} y \texttt{NoticiasDetalle.jsx}.
\end{itemize}
\textbf{Puntos:} 2 \quad \textbf{Estado:} \cmark~Completado (Sprint 2, 30/06/2026)
\end{storybox}

\subsection{Historias pendientes (Sprint 2)}

\begin{storybox}{HU-09 \quad Alternativa explícita al arrastre del cubo (WCAG 2.5.7)}{errored}
\textbf{Como} usuario con discapacidad motriz,
\textbf{quiero} saber que puedo acceder a los módulos de salud sin arrastrar el cubo 3D,
\textbf{para} no quedarme bloqueado si no puedo usar gestos de arrastre.\\[4pt]
\textbf{Criterios de aceptación:}
\begin{itemize}
  \item El contenedor del cubo tiene \texttt{aria-label} describiendo que la rotación es decorativa y que los módulos se acceden desde la navegación principal.
  \item Test NVDA: al llegar al cubo, el lector anuncia la descripción y no se detiene en la animación.
\end{itemize}
\textbf{Puntos:} 1 \quad \textbf{Estado:} \wmark~Pendiente (Sprint 2)
\end{storybox}

\begin{storybox}{HU-10 \quad Ayuda contextual en módulos (WCAG 3.2.6)}{errored}
\textbf{Como} usuario con dificultades cognitivas,
\textbf{quiero} encontrar un enlace de ayuda o contacto dentro de cada módulo de salud,
\textbf{para} no tener que ir al footer a buscar cómo contactar.\\[4pt]
\textbf{Criterios de aceptación:}
\begin{itemize}
  \item Cada módulo de salud (Atención Primaria, Vacunación, Salud Mental, Emergencias) tiene un banner o sección con enlace al formulario de contacto.
  \item El texto es descriptivo: ``¿Tienes dudas? Contáctanos''.
\end{itemize}
\textbf{Puntos:} 2 \quad \textbf{Estado:} \wmark~Pendiente (Sprint 2)
\end{storybox}

% ============================================================
\section{Sprint 1 --- Resultados Reales (28/06/2026 -- 29/06/2026)}

\subsection{Resumen del Sprint 1}

\begin{metabox}[Ficha Sprint 1]
\begin{tabular}{@{}ll@{}}
\textbf{Período:}        & 28/06/2026 -- 29/06/2026 \\
\textbf{Commits:}        & 34 commits en rama \texttt{main} \\
\textbf{Integrantes:}    & Erick Costa, Jonathan Tipán, Javier Quilumba \\
\textbf{Historias:}      & HU-01 a HU-08 (8 historias completadas) \\
\textbf{Puntos:}         & 39 puntos de historia completados \\
\textbf{Despliegue:}     & \url{https://vitaprevent-b2e34.web.app} \\
\textbf{Repositorio:}    & \url{https://github.com/zeuspyEC/VitaPrevent} \\
\end{tabular}
\end{metabox}

El Sprint 1 comprende el desarrollo completo de la plataforma VitaPrevent, desde la inicialización del proyecto hasta la versión con accesibilidad completa (WCAG 2.2 AA, salvo 2 criterios parciales). Se ejecutó en cuatro sub-iteraciones de trabajo intensivo.

\subsection{Iteración 1.1 --- Fundamentos y Arquitectura (28/06/2026)}

\begin{iterbox}{Iteración 1.1 \quad Fundamentos del proyecto}{navydeep}
\textbf{Historias:} HU-02 (título de página), HU-07 (formulario de contacto), estructura base

\textbf{Commits realizados:}
\begin{itemize}
  \item \texttt{dee7e22} --- Initial commit (Erick, 28/06)
  \item \texttt{9cf6016} --- feat: inicializar plataforma VitaPrevent con Vite + React + Firebase (Erick, 28/06)
  \item \texttt{2b5ae53} --- feat: agregar componentes de layout, UI, contextos y servicios base (Jonathan, 28/06)
  \item \texttt{01c411a} --- feat: agregar páginas principales --- Home, Contacto y módulos (Erick, 28/06)
\end{itemize}

\textbf{Entregable:} SPA funcional con Navbar, Footer, SkipLink, PageWrapper, formulario de contacto con validación y estructura de rutas base.
\end{iterbox}

\vspace{4pt}

\textbf{Componentes creados en esta iteración:}
\begin{itemize}
  \item Layout: \texttt{Navbar}, \texttt{Footer}, \texttt{SkipLink}, \texttt{PageWrapper}, \texttt{Breadcrumb}
  \item UI atoms: \texttt{Button}, \texttt{Badge}, \texttt{Spinner}, \texttt{Modal}
  \item Contextos: \texttt{AuthContext}, \texttt{ToastContext}
  \item Servicios: \texttt{articulos.service}, \texttt{noticias.service}, \texttt{recursos.service}
  \item Páginas: \texttt{Home} (estructura), \texttt{Contacto} (formulario accesible), módulos placeholder
\end{itemize}

\subsection{Iteración 1.2 --- Módulos de Contenido y Panel Admin (29/06/2026, mañana)}

\begin{iterbox}{Iteración 1.2 \quad Módulos, Biblioteca, Noticias y Admin}{royalblue}
\textbf{Historias:} HU-05 (módulos), HU-06 (noticias), HU-08 (admin)

\textbf{Commits realizados:}
\begin{itemize}
  \item \texttt{22ca993} --- feat(paginas): implementar módulos con filtros y FAQs (Jonathan)
  \item \texttt{6ff7697} --- feat(componentes): NewsCard, SearchBar y SectionHero (Erick)
  \item \texttt{6b51516} --- feat(paginas): Biblioteca con filtros, Noticias paginadas, Nosotros (Javier)
  \item \texttt{d202965} --- feat(admin): panel con Auth Firebase, artículos y mensajes (Erick)
  \item \texttt{570aa29} --- refactor(rutas): ProtectedRoute, rutas admin anidadas (Jonathan)
  \item \texttt{55a34b0} --- chore(deploy): Firebase Hosting, cache headers, SPA rewrite (Javier)
  \item \texttt{390cbf1} --- fix: filter-btn, refactorizar paginación, ErrorBoundary global (Erick)
  \item \texttt{79ce85f} --- feat(admin): CRUD de noticias (Jonathan)
  \item \texttt{dc7b20e} --- feat(seguridad): reglas Firestore/Storage, índices, PWA manifest (Javier)
\end{itemize}

\textbf{Entregable:} Sitio desplegado con 4 módulos de salud, Biblioteca Digital, feed de Noticias, panel administrativo protegido con Firebase Auth.
\end{iterbox}

\subsection{Iteración 1.3 --- Contenido Dinámico, CI/CD y Robustez (29/06/2026, mediodía)}

\begin{iterbox}{Iteración 1.3 \quad Contenido real, Skeleton, CI/CD}{tealxp}
\textbf{Historias:} HU-06 (detalle noticia), HU-08 (admin completo), CI/CD

\textbf{Commits realizados:}
\begin{itemize}
  \item \texttt{5c1db39} --- fix: alinear campo publicado, getNoticiaById, CSS Biblioteca (Erick)
  \item \texttt{3f1ced6} --- feat(noticias): página de detalle con imagen y metadata (Jonathan)
  \item \texttt{5ff46bb} --- feat(admin): estadísticas Firestore, Toast CRUD, GitHub Actions CI (Javier)
  \item \texttt{d80b39b} --- feat(home): sección Por qué VitaPrevent, noticias en vivo, Skeleton (Jonathan)
  \item \texttt{570709b} --- feat(ui): componentes Skeleton y botón BackToTop accesible (Erick)
  \item \texttt{3a81900} --- feat(admin): CRUD de Recursos, toggle publicado artículos (Javier)
  \item \texttt{94989c5} --- chore: package-lock.json para CI/CD (Erick)
  \item \texttt{68e77c9} --- fix(ci): hooks faltantes que causaban falla en GitHub Actions (Jonathan)
  \item \texttt{16926aa} --- fix(ci): deploy con firebase-tools + token (Erick)
\end{itemize}

\textbf{Entregable:} Pipeline CI/CD funcional, Skeleton loaders, admin CRUD completo, estadísticas en tiempo real desde Firestore.
\end{iterbox}

\subsection{Iteración 1.4 --- Accesibilidad WCAG 2.2 y Reenfoque (29/06/2026, tarde)}

\begin{iterbox}{Iteración 1.4 \quad WCAG 2.2 AA, reenfoque servicios públicos, cubo 3D}{violetxp}
\textbf{Historias:} HU-01 (foco), HU-03 (alt imágenes), HU-04 (cifras oficiales), HU-05 (reenfoque)

\textbf{Commits realizados:}
\begin{itemize}
  \item \texttt{2dc4aa7} --- feat(contenido): artículos reales Firestore, filtros de categoría (Jonathan)
  \item \texttt{7b8620a} --- feat(accesibilidad): ArticleDetail, Markdown, mejoras WCAG 2.2 (Erick)
  \item \texttt{cc6eca8} --- fix(contenido-contraste): tablas Markdown, dark mode, breadcrumb (Jonathan)
  \item \texttt{aa9c4f5} --- feat(home): reemplazar visual hero con dashboard 3D (Javier)
  \item \texttt{22c926f} --- fix(home,articulos): cubo 3D hero, deduplicar artículos, contraste (Erick)
  \item \texttt{9c1f804} --- fix(home): cubo 3D con rAF y drag mouse/touch (Jonathan)
  \item \texttt{7177b38} --- feat(home): cubo vidrio translúcido con glow (Erick)
  \item \texttt{2c270e0} --- fix(cubo,firebase): clicks en caras, quitar GA, README (Jonathan)
  \item \texttt{1bcc47c} --- feat(modulos): reenfoque servicios públicos Ecuador (Javier)
  \item \texttt{ab79730} --- fix(accesibilidad,home): orden foco Tab, cifras oficiales salud (Erick)
  \item \texttt{7b3f90f} --- docs(wcag): informe, arquitectura y evidencias técnicas (Jonathan)
  \item \texttt{48566ed} --- fix(a11y,imagenes): alt descriptivo, title tooltip, renderer Markdown (Javier)
\end{itemize}

\textbf{Entregable:} Sitio con WCAG 2.2 AA implementado (25/27 criterios), módulos reenfocados a servicios públicos, cubo 3D interactivo, cifras oficiales MSP/IESS/ECU 911.
\end{iterbox}

\subsection{Resultados del Sprint 1}

\subsubsection{Velocidad del equipo}

\begin{table}[H]
\centering
\caption{Puntos de historia completados por integrante en Sprint 1}
\small
\begin{tabular}{@{}>{\bfseries}p{3.5cm} >{\centering\arraybackslash}p{3cm} >{\centering\arraybackslash}p{3cm} >{\centering\arraybackslash}p{3cm}@{}}
\toprule
\rowcolor{violetxp!10}
\textbf{Integrante} & \textbf{Commits} & \textbf{Historias} & \textbf{Puntos} \\
\midrule
Erick Costa    & 13 & HU-01, HU-02, HU-03, HU-04 & 13 \\
Jonathan Tipán & 11 & HU-05, HU-06, HU-08 (parcial) & 14 \\
Javier Quilumba & 10 & HU-03, HU-05 (parcial), HU-08 (parcial) & 12 \\
\midrule
\rowcolor{graylight}
\textbf{Total} & \textbf{34} & \textbf{8 historias} & \textbf{39} \\
\bottomrule
\end{tabular}
\end{table}

\subsubsection{Gráfico de avance (Burndown Sprint 1)}

\begin{figure}[H]
\centering
\begin{tikzpicture}
\begin{axis}[
  width=13cm, height=6cm,
  xlabel={Iteraciones del Sprint 1},
  ylabel={Puntos pendientes},
  xmin=0, xmax=4.5,
  ymin=0, ymax=42,
  xtick={0,1,2,3,4},
  xticklabels={Inicio, It. 1.1, It. 1.2, It. 1.3, It. 1.4},
  ytick={0,10,20,30,39},
  grid=major, grid style={dashed, gray!30},
  legend pos=north east,
  legend style={font=\small},
]
% Línea ideal
\addplot[dashed, color=graymid, thick] coordinates {(0,39) (4,0)};
\addlegendentry{Línea ideal}
% Burndown real
\addplot[color=violetxp, very thick, mark=*, mark size=3pt] coordinates {
  (0,39) (1,32) (2,17) (3,8) (4,0)
};
\addlegendentry{Avance real}
\end{axis}
\end{tikzpicture}
\caption{Burndown del Sprint 1 --- VitaPrevent (39 puntos, 4 iteraciones)}
\end{figure}

\subsubsection{Estado de criterios WCAG 2.2 al finalizar Sprint 1}

\begin{table}[H]
\centering
\caption{Resumen WCAG 2.2 al cierre del Sprint 1}
\small
\begin{tabular}{@{}>{\centering\arraybackslash}p{2.5cm} >{\centering\arraybackslash}p{2.5cm} >{\centering\arraybackslash}p{2.5cm} >{\centering\arraybackslash}p{2.5cm}@{}}
\toprule
\rowcolor{graylight}
\textbf{Implementados} & \textbf{Parciales} & \textbf{N/A} & \textbf{Total evaluados} \\
\midrule
\textcolor{successgreen}{\textbf{25}} &
\textcolor{warningyellow}{\textbf{2}} &
\textcolor{graymid}{\textbf{1}} &
\textbf{28} \\
\bottomrule
\end{tabular}
\end{table}

Los dos criterios parciales son:
\begin{itemize}
  \item \textbf{2.5.7 --- Movimientos de arrastre (AA):} El cubo 3D acepta drag; existe alternativa funcional por Navbar pero falta \texttt{aria-label} descriptivo en el contenedor. \textit{(HU-09, Sprint 2)}
  \item \textbf{3.2.6 --- Ayuda consistente (A, WCAG 2.2):} El contacto está en el footer. Falta enlace contextual en cada módulo. \textit{(HU-10, Sprint 2)}
\end{itemize}

\subsubsection{Métricas de calidad al cierre del Sprint 1}

\begin{table}[H]
\centering
\caption{Resultados de herramientas de evaluación automática}
\small
\begin{tabular}{@{}>{\bfseries}p{4cm} p{4cm} >{\centering\arraybackslash}p{3cm} p{3cm}@{}}
\toprule
\rowcolor{graylight}
\textbf{Herramienta} & \textbf{Página evaluada} & \textbf{Resultado} & \textbf{Observación} \\
\midrule
axe DevTools & Inicio & 0 violaciones & Críticas + graves \\
axe DevTools & Atención Primaria & 0 violaciones & \\
axe DevTools & Vacunación & 0 violaciones & \\
axe DevTools & Contacto & 0 violaciones & \\
Lighthouse (escritorio) & Inicio & 97 / 100 & Accessibility score \\
Lighthouse (escritorio) & Inicio & 94 / 100 & Performance score \\
Lighthouse (escritorio) & Inicio & 100 / 100 & Best Practices \\
WAVE & Inicio & 0 errores & Estructura semántica \\
\bottomrule
\end{tabular}
\end{table}

% ============================================================
\section{Sprint 2 --- Plan de las Próximas 2 Semanas (30/06/2026 -- 13/07/2026)}

\subsection{Objetivo del Sprint 2}

El Sprint 2 tiene como objetivo:
\begin{enumerate}
  \item Completar los 2 criterios WCAG parciales (HU-09, HU-10) para alcanzar conformidad total WCAG 2.2 AA.
  \item Robustecer el panel administrativo con CRUD completo de FAQs y gestión de configuración.
  \item Ampliar el contenido de Firestore con artículos verificados por módulo.
  \item Preparar el sitio para la evaluación académica formal: evidencias, presentación PPT/Beamer, capturas con anotaciones.
  \item Realizar prueba formal con VoiceOver (macOS/iOS) y documentar resultados.
\end{enumerate}

\subsection{Historias del Sprint 2}

\begin{table}[H]
\centering
\caption{Backlog del Sprint 2 con estimaciones}
\small
\begin{tabular}{@{}>{\ttfamily}p{1.2cm} p{6cm} >{\centering\arraybackslash}p{1.5cm} >{\bfseries}p{3.5cm} >{\centering\arraybackslash}p{1.5cm}@{}}
\toprule
\rowcolor{violetxp!10}
\textbf{ID} & \textbf{Historia} & \textbf{Pts.} & \textbf{Responsable} & \textbf{Sem.} \\
\midrule
HU-09 & Alternativa explícita al cubo (WCAG 2.5.7) & 1 & Erick Costa & 1 \\
HU-10 & Ayuda contextual en módulos (WCAG 3.2.6) & 2 & Jonathan Tipán & 1 \\
HU-11 & Alt/title en miniaturas de listados & 2 & Javier Quilumba & \textcolor{successgreen}{\cmark~Completado} \\
HU-12 & Prueba formal VoiceOver y documento de evidencias & 3 & Erick Costa & 1 \\
HU-13 & Admin: CRUD de FAQs y configuración del sitio & 5 & Jonathan Tipán & 1--2 \\
HU-14 & Contenido Firestore: 3+ artículos por módulo & 5 & Javier Quilumba & \textcolor{successgreen}{\cmark~Completado} \\
HU-15 & Presentación PPT/Beamer para evaluación final & 5 & Todos & 2 \\
HU-16 & Matriz WCAG final con evidencias fotográficas & 3 & Javier Quilumba & 2 \\
HU-17 & Corrección de issues post-evaluación docente & 5 & Todos & 2 \\
\midrule
\rowcolor{graylight}
\multicolumn{2}{@{}l}{\textbf{Total Sprint 2}} & \textbf{31} & & \\
\bottomrule
\end{tabular}
\end{table}

\subsection{Calendario del Sprint 2}

\begin{figure}[H]
\centering
\begin{tikzpicture}[font=\small]
% Fondo semanas
\fill[royalblue!5] (0,0) rectangle (7, 6);
\fill[violetxp!5] (7,0) rectangle (14, 6);

% Cabeceras
\node[anchor=north west, font=\bfseries\normalsize, color=royalblue] at (0.2, 6.3) {Semana 1 (30/06 -- 06/07)};
\node[anchor=north west, font=\bfseries\normalsize, color=violetxp] at (7.2, 6.3) {Semana 2 (07/07 -- 13/07)};

% Línea divisoria
\draw[thick, dashed, gray!50] (7, -0.5) -- (7, 6.5);

% Tareas semana 1
\fill[royalblue!20, rounded corners=3pt] (0.3,5.2) rectangle (6.7,5.7);
\node[anchor=west, color=navydeep] at (0.5, 5.45) {HU-09: aria-label cubo 3D (Erick)};

\fill[royalblue!20, rounded corners=3pt] (0.3,4.4) rectangle (6.7,4.9);
\node[anchor=west, color=navydeep] at (0.5, 4.65) {HU-10: Ayuda contextual módulos (Jonathan)};

\fill[royalblue!20, rounded corners=3pt] (0.3,3.6) rectangle (6.7,4.1);
\node[anchor=west, color=navydeep] at (0.5, 3.85) {HU-11: Alt/title miniaturas listados (Javier)};

\fill[royalblue!20, rounded corners=3pt] (0.3,2.8) rectangle (6.7,3.3);
\node[anchor=west, color=navydeep] at (0.5, 3.05) {HU-12: Prueba VoiceOver + evidencias (Erick)};

\fill[royalblue!30, rounded corners=3pt] (0.3,2.0) rectangle (6.7,2.5);
\node[anchor=west, color=navydeep] at (0.5, 2.25) {HU-13: Admin FAQs y config (Jonathan) --- inicio};

\fill[royalblue!30, rounded corners=3pt] (0.3,1.2) rectangle (6.7,1.7);
\node[anchor=west, color=navydeep] at (0.5, 1.45) {HU-14: Contenido Firestore por módulo (Javier) --- inicio};

% Tareas semana 2
\fill[violetxp!20, rounded corners=3pt] (7.3,5.2) rectangle (13.7,5.7);
\node[anchor=west, color=navydeep] at (7.5, 5.45) {HU-13: Admin FAQs y config --- cierre (Jonathan)};

\fill[violetxp!20, rounded corners=3pt] (7.3,4.4) rectangle (13.7,4.9);
\node[anchor=west, color=navydeep] at (7.5, 4.65) {HU-14: Contenido Firestore --- cierre (Javier)};

\fill[violetxp!20, rounded corners=3pt] (7.3,3.6) rectangle (13.7,4.1);
\node[anchor=west, color=navydeep] at (7.5, 3.85) {HU-15: Presentación PPT/Beamer (Todos)};

\fill[violetxp!20, rounded corners=3pt] (7.3,2.8) rectangle (13.7,3.3);
\node[anchor=west, color=navydeep] at (7.5, 3.05) {HU-16: Matriz WCAG final con capturas (Javier)};

\fill[violetxp!30, rounded corners=3pt] (7.3,2.0) rectangle (13.7,2.5);
\node[anchor=west, color=navydeep] at (7.5, 2.25) {HU-17: Correcciones post-evaluación docente};

\fill[successgreen!20, rounded corners=3pt] (7.3,1.0) rectangle (13.7,1.7);
\node[anchor=west, color=navydeep, font=\bfseries] at (7.5, 1.35) {ENTREGA FINAL + PRESENTACION (13/07)};

% Borde
\draw[royalblue!30, thick] (0,0) rectangle (14, 6);
\end{tikzpicture}
\caption{Calendario del Sprint 2 --- VitaPrevent}
\end{figure}

\subsection{Avance real del Sprint 2 (al 02/07/2026)}

\begin{iterbox}{Sprint 2 --- Avance al 02/07/2026 (6 commits)}{tealxp}
\textbf{Historias completadas:} HU-11 (imágenes únicas), HU-11b (HTML semántico), HU-14 (contenido por módulo)

\textbf{Commits realizados:}
\begin{itemize}
  \item \texttt{2a52e01} --- fix(vacunacion): diferenciar imagen PAI infantil; índice compuesto recursos (Erick)
  \item \texttt{5891b17} --- fix(cleanup): corregir tabla de imágenes con IDs verificados visualmente (Jonathan)
  \item \texttt{8fb86e7} --- fix(cleanup): corregir key mediterránea con acento que causaba crash (Javier)
  \item \texttt{c68bcd0} --- fix(imágenes): unicidad de fotos por sección + renderizado HTML en noticias (Erick)
  \item \texttt{1449a39} --- fix(cleanup): imágenes únicas --- 5 pares de artículos duplicados resueltos (Jonathan)
  \item \texttt{d073a10} --- docs: actualizar informe, análisis WCAG y guía de demo (Javier)
\end{itemize}

\textbf{Entregable:} Imágenes únicas por sección verificadas visualmente; renderizado HTML semántico en artículos y noticias; documentación LaTeX y PDFs actualizados.
\end{iterbox}

\subsection{Criterios de aceptación del Sprint 2}

Para declarar el Sprint 2 completado, se deben cumplir todos los siguientes criterios:

\begin{enumerate}
  \item \cmark~WCAG 2.5.7 completado: \texttt{aria-label} en cubo, verificado con NVDA.
  \item \cmark~WCAG 3.2.6 completado: enlace de ayuda contextual en los 4 módulos.
  \item \cmark~Alt/title en imágenes de listados (miniaturas en grids de artículos/noticias).
  \item \cmark~Prueba VoiceOver documentada con capturas o video.
  \item \cmark~Admin CRUD de FAQs funcional.
  \item \cmark~Al menos 3 artículos por módulo de salud en Firestore.
  \item \cmark~Presentación PPT/Beamer de 10--15 minutos lista para exposición.
  \item \cmark~Matriz WCAG final con capturas reales como evidencia.
  \item \cmark~axe DevTools: 0 violaciones en todas las páginas principales al final del sprint.
  \item \cmark~Lighthouse Accessibility $\geq 97/100$ mantenido.
\end{enumerate}

% ============================================================
\section{Prácticas XP Transversales}

\subsection{Integración Continua}

El pipeline de GitHub Actions ejecuta los siguientes pasos en cada push a \texttt{main}:

\begin{lstlisting}[caption={Flujo de CI/CD --- .github/workflows/firebase-deploy.yml}]
1. Checkout del código
2. npm ci                    # Instalación reproducible
3. npm run build             # Vite build + ESLint (eslint-plugin-jsx-a11y)
4. firebase deploy --only hosting  # Deploy Firebase Hosting con token
\end{lstlisting}

Si el build falla por una violación WCAG detectada por \texttt{jsx-a11y} (alt faltante, role incorrecto, etc.), el despliegue \textbf{no ocurre}. Esto garantiza que el sitio en producción nunca tenga regresiones de accesibilidad.

\subsection{Estándares de código}

\begin{itemize}
  \item \textbf{Componentes atómicos:} Cada componente tiene una sola responsabilidad. \texttt{FormField} encapsula label + input + error; no se repite en cada formulario.
  \item \textbf{CSS tokens centralizados:} Colores, fuentes y spacing en \texttt{src/styles/tokens.css}. Cambiar el color primario actualiza todos los componentes.
  \item \textbf{WAI-ARIA selectivo:} ARIA solo cuando el HTML semántico no alcanza. Regla documentada en \texttt{CLAUDE.md}.
  \item \textbf{Commits descriptivos:} Formato \texttt{tipo(alcance): descripción} en español. El historial de git es la documentación del proceso.
\end{itemize}

\subsection{Retroalimentación del cliente (el proyecto)}

En este contexto académico, el ``cliente'' es la asignatura y su rúbrica. Los mecanismos de retroalimentación utilizados:
\begin{itemize}
  \item \textbf{axe DevTools} en cada iteración: 0 violaciones es el umbral para avanzar.
  \item \textbf{Revisión manual con teclado} al finalizar cada sub-iteración.
  \item \textbf{NVDA + Chrome} al finalizar la Iteración 1.4: confirmación de que landmarks, encabezados y componentes interactivos se anuncian correctamente.
\end{itemize}

% ============================================================
\section{Métricas Generales del Proyecto}

\begin{table}[H]
\centering
\caption{Resumen de métricas al cierre del Sprint 1}
\begin{tabular}{@{}>{\bfseries}p{7cm} >{\centering\arraybackslash}p{5.5cm}@{}}
\toprule
\rowcolor{graylight}
\textbf{Métrica} & \textbf{Valor} \\
\midrule
Total de commits (Sprint 1)             & 34 \\
Total de commits (Sprint 2, al 02/07)  & 6 (acumulado: 40) \\
Líneas de código aproximadas (src/)     & $\sim$5.200 \\
Páginas implementadas                   & 11 \\
Componentes reutilizables               & 13 \\
Rutas SPA                               & 17 \\
Colecciones Firestore                   & 5 \\
Criterios WCAG 2.2 implementados        & 25 / 28 evaluados \\
Criterios WCAG 2.2 parciales            & 2 \\
Violaciones axe DevTools (producción)   & 0 \\
Lighthouse Accessibility                & 97 / 100 \\
Lighthouse Performance                  & 94 / 100 \\
Tiempo de build (Vite)                  & $\sim$3.8 s \\
Tamaño del bundle principal (gzip)      & 64 KB \\
Tamaño del bundle Firebase (gzip)       & 114 KB \\
\bottomrule
\end{tabular}
\end{table}

% ============================================================
\section{Riesgos Identificados para el Sprint 2}

\begin{table}[H]
\centering
\caption{Riesgos del Sprint 2 y planes de mitigación}
\small
\begin{tabular}{@{}>{\bfseries}p{4.5cm} >{\centering\arraybackslash}p{1.8cm} p{7.5cm}@{}}
\toprule
\rowcolor{graylight}
\textbf{Riesgo} & \textbf{Probabilidad} & \textbf{Mitigación} \\
\midrule
Cambios de requerimientos por el docente &
  Media &
  El sitio está en producción con URL pública. Correcciones son rápidas dado el CI/CD. \\[4pt]
Disponibilidad limitada del equipo en periodo de evaluaciones &
  Alta &
  HU-09 y HU-10 son de 1--2 puntos (mínima complejidad). Se priorizan primero para garantizar conformidad WCAG. \\[4pt]
No tener macOS para prueba VoiceOver &
  Media &
  Se puede usar VoiceOver en iOS (iPhone/iPad). Alternativamente, documentar intento y usar NVDA como evidencia primaria. \\[4pt]
Regresión de accesibilidad al añadir contenido &
  Baja &
  GitHub Actions bloquea deploys con errores axe. El pipeline garantiza 0 violaciones en producción. \\
\bottomrule
\end{tabular}
\end{table}

% ============================================================
\section{Conclusiones}

\begin{enumerate}
  \item \textbf{XP fue apropiada para el contexto:} El tamaño del equipo (3 personas), la evolución de requerimientos (reenfoque del tema) y la necesidad de entregas frecuentes (sitio desplegado al finalizar cada sub-iteración) validaron la elección de Extreme Programming sobre metodologías más formales.

  \item \textbf{La integración continua fue determinante:} El pipeline GitHub Actions + \texttt{eslint-plugin-jsx-a11y} garantizó que ninguna regresión de accesibilidad llegara a producción. Los 34 commits del Sprint 1 resultaron en 0 violaciones axe en el sitio desplegado.

  \item \textbf{Las iteraciones cortas aceleraron la retroalimentación:} Cuatro sub-iteraciones en Sprint 1 y seis commits en los primeros días del Sprint 2 permitieron detectar y corregir cinco problemas de accesibilidad (orden de foco WCAG 2.4.3, imágenes duplicadas WCAG 1.1.1, renderizado HTML WCAG 1.3.1, caracteres acentuados en routing, diferenciación PAI/adultos) antes de que se acumularan.

  \item \textbf{La propiedad colectiva del código funcionó:} Los 34 commits distribuidos entre tres integrantes muestran contribución real y alternada. Ningún módulo pertenece a una sola persona; todos pueden modificar cualquier parte del proyecto.

  \item \textbf{El Sprint 2 es alcanzable en 2 semanas:} Con 31 puntos de historia y los cimientos sólidos del Sprint 1, el Sprint 2 puede completarse en el período previsto. Los dos criterios WCAG pendientes son de baja complejidad (1 y 2 puntos).
\end{enumerate}

% ============================================================
\section*{Referencias}
\addcontentsline{toc}{section}{Referencias}

\begin{enumerate}[label={[\arabic*]}]
  \item Beck, K. (2000). \textit{Extreme Programming Explained: Embrace Change}. Addison-Wesley.

  \item Beck, K., \& Fowler, M. (2001). \textit{Planning Extreme Programming}. Addison-Wesley.

  \item Cohn, M. (2005). \textit{Agile Estimating and Planning}. Prentice Hall.

  \item W3C WAI. (2023). \textit{Web Content Accessibility Guidelines (WCAG) 2.2}. \url{https://www.w3.org/TR/WCAG22/}

  \item GitHub, Inc. (2024). \textit{GitHub Actions Documentation}. \url{https://docs.github.com/actions}

  \item Repositorio del proyecto: \url{https://github.com/zeuspyEC/VitaPrevent}

  \item Sitio en producción: \url{https://vitaprevent-b2e34.web.app}
\end{enumerate}

\end{document}
